%----------
%	INSULTOS - LLMs
%----------	


Similarly to the approaches detailed in Sections~\ref{ageneral-genai} and ~\ref{cnegativo-genai} for ML and DL, the last classification approach that was employed for the "Insultos" task was using generative models and zero-shot learning. In total, there were 64 total experiments for the two types of contexts and 4 different prompts, as defined in Section~\ref{implementation_genai}, defining a total of 8 different windows that can be observed in Table~\ref{tab:window-metrics-insultos-genai}. In particular, W1-W4 represented experiments using "Context\_0", while W5-W8 used "Context\_1." Within each context, the first window (W1 and W5) utilized "Prompt\_0\_EN", the second (W2 and W6) employed "Prompt\_0\_ES", the third (W3 and W7) applied "Prompt\_1\_EN", and the final window (W4 and W8) used "Prompt\_1\_ES".


\begin{table}[H]
    \centering
    \begingroup
    \small % Change to desired font size
    \ttabbox[\FBwidth]{
        \caption{INS Average Evaluation Metrics per Window in GenAI Experiments}
        \label{tab:window-metrics-insultos-genai} % Unique label for the table
    }{
        \begin{tabularx}{\textwidth}{
            X
            X
            X
            >{\centering\arraybackslash}X
            >{\centering\arraybackslash}X
            >{\columncolor{gray!15}\centering\arraybackslash}X % Highlight this column
            >{\centering\arraybackslash}X
            >{\centering\arraybackslash}X
            >{\centering\arraybackslash}X
        }
            \hline
            \rowcolor{gray!30} % Adding a gray color to the header row
            \textbf{} & \textbf{Context} & \textbf{Prompt} & \textbf{Acc. (Std.)} & \textbf{Time (s)} & \textbf{W. Avg. F1-Score} & \textbf{Sex./Mis. F1-Score} & \textbf{Gen. F1-Score} & \textbf{D. Dañar F1-Score} \\
            \hline
            W1 & 0 & 0\_EN & 0.487 (0.088) & 85.778 & 0.589 & 0.175 & 0.645 & 0.108 \\
            \hline
            W2 & 0 & 0\_ES & 0.466 (0.098) & 97.133 & 0.520 & 0.331 & 0.610 & 0.140 \\
            \hline
            W3 & 0 & 1\_EN & 0.466 (0.103) & 92.310 & 0.557 & 0.167 & 0.629 & 0.112 \\
            \hline
            W4 & 0 & 1\_ES & 0.480 (0.130) & 89.707 & 0.528 & 0.372 & 0.632 & 0.139 \\
            \hline
            W5 & 1 & 0\_EN & 0.501 (0.061) & 95.186 & 0.591 & 0.321 & 0.654 & 0.069 \\
            \hline
            W6 & 1 & 0\_ES & 0.434 (0.128) & 88.408 & 0.478 & 0.362 & 0.579 & 0.092 \\
            \hline
            W7 & 1 & 1\_EN & 0.476 (0.110) & 94.778 & 0.547 & 0.222 & 0.640 & 0.153 \\
            \hline
            W8 & 1 & 1\_ES & 0.465 (0.136) & 93.928 & 0.531 & 0.330 & 0.624 & 0.070 \\
            \hline
            \multicolumn{9}{l}{Source: 'Training\_GenAI.xlsx'} \\
            \hline
        \end{tabularx}
    }
    \endgroup
\end{table}


In Figure~\ref{fig:exp_insultos_genai} we can observe the distribution of experiments with the different windows. Overall, we could see a decreasing trend in the weighted average F1-Score, although the first windows of each Context (W1 and W5) were the best-performing windows with 0.589 and 0.591. However, when comparing the different contexts we saw that there was not a significant difference in performance with a mean weighted average F1-Score of 0.548 for Context 0 and 0.537 for Context 1. Additionally, we observed that odd-numbered windows, corresponding to prompts in English, outperformed those in Spanish. These findings might suggest that while the contexts may not have significantly influenced the classification, the prompts did, with "Prompt\_0\_EN" obtaining the highest performance.

On the other hand, Figure~\ref{fig:exp_insultos_classes_genai} shows the distribution of experiments across different classes. We observed a consistent and stable trend for the "Genéricos" and "Deseo de Dañar" classes, with the highest mean F1-Scores achieved in W5 and W7, respectively, under different contexts and prompt settings. However, in "Deseo de Dañar" we could see a slightly decreasing trend, except for W7 (Context 1, Prompt\_1\_EN), where the highest F1-Score was achieved with 0.153. Additionally, we could see that the best overall performing window was W4 with F1-Scores of 0.372, 0.632 and 0.139 for "Sexistas/Misóginos", "Genéricos" and "Deseo de Dañar", respectively.


\begin{figure}[H]
	\ffigbox[\FBwidth]
	{\caption{Distribution of the GenAI Experiments over Time for "Insultos" with the Different Windows and Weighted Avg. F1-Score Trend Line}\label{fig:exp_insultos_genai}}
 % [Aquí ponemos como queremos que aparezca en el índice, sin la referencia]{Aquí como queremos que aparezca en el documento, con la referencia}
	{\includegraphics[scale=0.25]{Plantilla_TFG_ingles_2019/imagenes/GenAI_Experiments in _Insultos_ weighted.pdf}}
\end{figure}


\begin{figure}[H]
	\ffigbox[\FBwidth]
	{\caption{Distribution of the GenAI Experiments over Time for "Insultos", per Class with the Different Windows, and Trend Lines for Each Class}\label{fig:exp_insultos_classes_genai}}
 % [Aquí ponemos como queremos que aparezca en el índice, sin la referencia]{Aquí como queremos que aparezca en el documento, con la referencia}
	{\includegraphics[scale=0.25]{Plantilla_TFG_ingles_2019/imagenes/GenAI_Experiments in _Insultos_f1score.pdf}}
\end{figure}


\newpage

A summary of the top 5 best-performing generative models for this task is presented in Table~\ref{tab:genai-cnegativo-bestmodels_metrics}. The results showed a consistent weighted average F1-Score of around 0.6, with accuracy ranging between 0.49 and 0.51. The "Genéricos" class achieved the highest F1-Score of 0.669 and demonstrated the most consistent performance, likely due to its generic nature, which may align more with the pre-trained models' general understanding. On the other hand, the "Deseo de Dañar" and "Sexistas/Misóginos" classes exhibited more variability, with the highest F1-Score for "Deseo de Dañar" reaching 0.600, closely following "Genéricos." However, the model struggled with the "Sexistas/Misóginos" class, achieving a maximum F1-Score of 0.347.


\begin{table}[H]
    \centering
    \begingroup
    \small % Change to desired font size
    \ttabbox[\FBwidth]{
        \caption{INS Evaluation Metrics for the Best GenAI Models}
        \label{tab:genai-cnegativo-bestmodels_metrics} % Unique label for the table
    }{
        \begin{tabularx}{\textwidth}{
            X
            >{\centering\arraybackslash}X
            >{\columncolor{gray!15}\centering\arraybackslash}X % Highlight this column
            >{\centering\arraybackslash}X
            >{\centering\arraybackslash}X
            >{\centering\arraybackslash}X
            >{\centering\arraybackslash}X
        }
            \hline
            \rowcolor{gray!30} % Adding a gray color to the header row
            \textbf{} & \textbf{Acc. (Std.)} & \textbf{Weighted Avg. F1-Score} & \textbf{Sex./Mis. F1-Score} & \textbf{Genéricos F1-Score} & \textbf{D. Dañar F1-Score} \\
            \hline
            \multicolumn{6}{l}{\textbf{Naive Models}} \\
            \hline
            \insmajor &  0.485 (0.035) & 0.318 & 0.000 & 0.653 & 0.000 \\
            \hline
            \insrandom & 0.369 (0.033) & 0.369 & 0.231 & 0.484 & 0.279 \\
            \hline
            \multicolumn{6}{l}{\textbf{Zero-Shot Models}} \\
            \hline
            \insllamathreeonea & 0.490 (0.110) & 0.607 & 0.133 & 0.651 & 0.600 \\
            \hline
            \insgemmatwoba & 0.490 (0.079) & 0.606 & 0.133 & 0.653 & 0.500 \\
            \hline
            \insgemmatwobb & 0.515 (0.105) & 0.606 & 0.347 & 0.669 & 0.286 \\
            \hline
            \insllamathreeoneb & 0.510 (0.030) & 0.603 & 0.338 & 0.659 & 0.500  \\
            \hline
            \insllamathreeonec & 0.510 (0.049) & 0.602 & 0.347 & 0.657 & 0.400 \\
            \hline
            \multicolumn{6}{l}{Source: 'Training\_GenAI.xlsx'} \\
            \hline
        \end{tabularx}
    }
    \endgroup
\end{table}


After these findings, the model \textbf{\insllamathreeonea} was selected as the best-performing model for the Generative AI approach in the "Insultos" task, as it obtained the highest F1-Scores for weighted average, "Deseo de Dañar" and "Genéricos" with 0.607, 0.600 and 0.651, and decent performance for "Sexistas/Misóginos" (0.133). This model generally outperformed the two naive baseline models, except in the "Sexistas/Misóginos" class. In Figure~\ref{fig:exp_insultos_genai_bestmodel} we can see in more detail the confusion matrix for model \textbf{\insllamathreeonea}. Notably, the model correctly classified 94 instances of "Genéricos" but struggled with misclassifying this class as "Deseo de Dañar" (52 instances) and "Sexistas/Misóginos" (43 instances). This may suggest that the zero-shot approach produced a solid performance of the "Genéricos" class, but some refinement is needed on more specific classes, such as "Sexistas/Misóginos" and "Deseo de Dañar".


% Best Model
\begin{figure}[H]
	\ffigbox[\FBwidth]
	{\caption{Confusion Matrix for the Best GenAI Model (\textbf{\insllamathreeonea}) for "Insultos"}\label{fig:exp_insultos_genai_bestmodel}}
 % [Aquí ponemos como queremos que aparezca en el índice, sin la referencia]{Aquí como queremos que aparezca en el documento, con la referencia}
	{\includegraphics[scale=0.15]{Plantilla_TFG_ingles_2019/imagenes/GenAI_Experiments in _Insultos_ best_model.pdf}}
\end{figure}

\newpage